\problem{2.46}{Stationary quantum states on a ring}
A ring of circumference $L$ imposes periodicity
$\psi(x+L)=\psi(x)$. Free solutions are
$\ee^{\ii kx}$; periodicity requires
$k_nL=2\pi n$ for integer $n$.
Normalization over one circumference gives
\begin{equation}
\boxed{\psi_n(x)=\frac1{\sqrt L}\ee^{2\pi\ii nx/L}},
\qquad
\boxed{E_n=\frac{2\pi^2\hbar^2n^2}{mL^2}},
\quad n\in\mathbb Z.
\end{equation}
For each positive $n$, the independent
$\psi_n$ and $\psi_{-n}$ carry opposite currents
and have equal energy. The $n=0$ state is
nondegenerate. The Wronskian theorem from
Problem 2.44 assumes wavefunctions decay
at spatial infinity on a connected line;
a circle is periodic rather than a line
with vanishing tails, so its constant
Wronskian need not be zero.

\problem{2.47}{Qualitative behavior of a double finite well}
At zero separation, the two attractive wells merge
into one wider region, and the lowest wavefunction
is nodeless and broadly distributed; the first excited
function has one antisymmetric node. At intermediate
separation the lowest pair becomes a symmetric and
antisymmetric combination of wavefunctions localized
near the left and right wells. At very large separation,
the overlap between the two centers becomes tiny:
the even and odd combinations have nearly the
same energy and become effectively independent.
The ground energy tends toward the isolated-well
ground energy from below, whereas the antisymmetric
partner tends toward it from above.
The splitting shrinks as the barrier separation
increases. In this model the electron favors closer
nuclei because the symmetric bonding state's
energy falls as the wells overlap. Real molecules
also involve internuclear repulsion, electron--electron
effects and three-dimensional geometry, so the
simple potential is only an analogy.
\check{This is deliberately a qualitative answer:
the textbook specifies no calculation for this problem.}

\problem{2.48}{Bound states beside a hard wall}
The potential is infinite for $x<0$, equals
$-V_0$ for $0<x<a$, and is zero for $x>a$,
where $V_0=32\hbar^2/(ma^2)$. For bound
states $-V_0<E<0$, set
\begin{equation}
z=ka=\frac{a\sqrt{2m(E+V_0)}}{\hbar},
\quad q=\kappa a=\frac{a\sqrt{-2mE}}{\hbar},
\quad z^2+q^2=64.
\end{equation}
The wall forces $\psi=C\sin(kx)$ inside;
outside, $\psi=C\sin z\,\ee^{-\kappa(x-a)}$.
The derivative matching gives
\begin{equation}
\boxed{-z\cot z=\sqrt{64-z^2}},\quad0<z<8.
\end{equation}
The three roots of this equation occur at
\begin{equation}
z\simeq2.785902,\quad5.521446,\quad7.957321,
\end{equation}
so there are three bound states. Their energies are
$E=-\hbar^2q^2/(2ma^2)$ with
$q=\sqrt{64-z^2}$.
For the highest state, $q\simeq0.825248$.
Integrating its density on either side of $a$,
\begin{align}
I_{\mathrm{in}}&=C^2a\left(
\frac12-\frac{\sin(2z)}{4z}\right),\\
I_{\mathrm{out}}&=C^2a\frac{\sin^2z}{2q},\\
P_{\mathrm{out}}&=\frac{I_{\mathrm{out}}}{I_{\mathrm{in}}+
I_{\mathrm{out}}}
=\boxed{0.542041\ \text{(approximately)}}.
\end{align}
\check{Despite $E<0$, the shallowest bound state
has more probability outside the finite well
than inside it.}

\problem{2.49}{Exact coherent oscillator wave packet}
Write $\ell^2=\hbar/(m\omega)$ and
$X(t)=x_0\cos\omega t$. The exact packet
given in the textbook can be written
\begin{align}
\Psi(x,t)&=(\pi\ell^2)^{-1/4}
\exp\Bigg[-\frac1{2\ell^2}
\left(x^2+\frac{x_0^2}2
(1+\ee^{-2\ii\omega t})
-2x_0x\ee^{-\ii\omega t}\right)
-\frac{\ii\omega t}{2}\Bigg].
\end{align}
To check it, denote its exponent by $S$ and use
$\partial_x\Psi=S_x\Psi$,
$\partial_x^2\Psi=(S_{xx}+S_x^2)\Psi$,
$\partial_t\Psi=S_t\Psi$. Inserting
\begin{align}
S_x&=-x/\ell^2+x_0\ee^{-\ii\omega t}/\ell^2,\\
S_{xx}&=-1/\ell^2,\\
S_t&=\frac{\ii\omega x_0^2}{\ell^2}\ee^{-2\ii\omega t}
-\frac{\ii\omega x_0 x}{\ell^2}\ee^{-\ii\omega t}
-\frac{\ii\omega}{2}
\end{align}
into $i\hbar S_t=
-\hbar^2(S_{xx}+S_x^2)/(2m)+m\omega^2x^2/2$
verifies the Schrödinger equation term by term.
Taking the real part of $S$,
\begin{equation}
\boxed{|\Psi(x,t)|^2=
\frac1{\sqrt{\pi}\ell}
\exp\left[-\frac{(x-x_0\cos\omega t)^2}{\ell^2}\right]}.
\end{equation}
The width remains constant:
$\sigma_x=\ell/\sqrt2$.
Gaussian integration gives
$\boxed{\avg x=x_0\cos\omega t}$ and
$\boxed{\avg p=-m\omega x_0\sin\omega t}$.
Then $\dd\avg p/\dd t=-m\omega^2\avg x$,
precisely Ehrenfest's theorem.

\problem{2.50}{A translating attractive delta well}
Let $V=-\alpha\delta(x-vt)$ and
$\kappa=m\alpha/\hbar^2$.
Start with the stationary bound state
$\psi_0(u)=\sqrt\kappa\ee^{-\kappa|u|}$
at energy $E_0=-m\alpha^2/(2\hbar^2)$.
A Galilean boost of velocity $v$ gives
\begin{equation}
\boxed{\Psi(x,t)=
\sqrt\kappa\ee^{-\kappa|x-vt|}
\exp\left[\frac{\ii}{\hbar}
\left(mvx-\left(E_0+\frac12mv^2\right)t\right)\right]}.
\end{equation}
For $u=x-vt$, differentiation of the phase
produces a kinetic-energy shift $mv^2/2$
and cross terms that cancel the translation
term $-v\partial_u\psi_0$. The remaining
equation is exactly the stationary delta-well
equation, including the distributional cusp.
Therefore this is a normalized exact solution.
Its instantaneous Hamiltonian expectation is
\begin{equation}
\boxed{\avg H=E_0+\frac12mv^2}.
\end{equation}
The center of the density moves at $vt$ while
the internal binding energy remains $E_0$.

\problem{2.51}{Gaussian free fall in a constant force}
Let $\Psi_0(y,t)$ be a freely evolving Gaussian
with $\avg y_0=0$ and zero mean momentum.
Define $y=x+gt^2/2$ and
\begin{equation}
\boxed{\Psi(x,t)=
\Psi_0\left(x+\frac12gt^2,t\right)
\exp\left[-\frac{\ii mgt}{\hbar}
\left(x+\frac16gt^2\right)\right]}.
\end{equation}
If $\Phi$ denotes the phase, its derivatives are
$\partial_x\Phi=-mgt/\hbar$ and
$\partial_t\Phi=-mg(x+gt^2/2)/\hbar$.
The chain rule gives
$\partial_t\Psi_0(y,t)=
(\partial_t\Psi_0)_y+gt\,\partial_y\Psi_0$;
the extra convective derivative cancels the
cross term from differentiating the phase
in the kinetic operator. The remaining term
is the required gravitational potential $mgx$.
Since the phase has unit modulus, the density
is simply translated:
\begin{equation}
|\Psi(x,t)|^2=
|\Psi_0(x+gt^2/2,t)|^2,\qquad
\boxed{\avg x(t)=-\frac12gt^2}.
\end{equation}
The center follows classical free fall although
the Gaussian simultaneously spreads.

\problem{2.52}{A reflectionless hyperbolic-secant well}
Let
\begin{equation}
V(x)=-\frac{\hbar^2a^2}{m}\operatorname{sech}^2(ax).
\end{equation}
The potential is even, approaches zero at both
infinities, and reaches $-\hbar^2a^2/m$ at zero.
For $\psi_0=A\operatorname{sech}(ax)$,
the identity
\begin{equation}
\frac{\psi_0''}{\psi_0}=
a^2[1-2\operatorname{sech}^2(ax)]
\end{equation}
shows
$\hat H\psi_0=-\hbar^2a^2\psi_0/(2m)$.
As $\int_{-\infty}^{\infty}
\operatorname{sech}^2(ax)\dd x=2/a$,
\begin{equation}
\boxed{\psi_0=\sqrt{\frac a2}\operatorname{sech}(ax),
\qquad E_0=-\frac{\hbar^2a^2}{2m}}.
\end{equation}
For $E=\hbar^2k^2/(2m)>0$,
take the exact scattering solution
\begin{equation}
\psi_k(x)=A\frac{\ii k-a\tanh(ax)}{\ii k+a}
\ee^{\ii kx}.
\end{equation}
Using $(\tanh ax)'=a\operatorname{sech}^2(ax)$
and $(\operatorname{sech}^2 ax)'=
-2a\operatorname{sech}^2(ax)\tanh(ax)$
in $\psi_k''$ verifies $\hat H\psi_k=E\psi_k$.
At $x\to-\infty$, $\tanh(ax)\to-1$,
so $\psi_k\to A\ee^{\ii kx}$, with no
left-moving reflected component. At
$x\to+\infty$,
\begin{equation}
\psi_k\to A
\frac{\ii k-a}{\ii k+a}\ee^{\ii kx}.
\end{equation}
The amplitude ratio has magnitude one, hence
\begin{equation}
\boxed{R=0,\qquad T=1}
\end{equation}
at every positive incident energy: an exactly
reflectionless potential.

\problem{2.53}{Scattering matrices for symmetric localized wells}
For a real localized potential, let $A,G$ be
the incoming amplitudes from left and right,
and $B,F$ the outgoing amplitudes:
\begin{equation}
\begin{pmatrix}B\\F\end{pmatrix}
=S\begin{pmatrix}A\\G\end{pmatrix}.
\end{equation}
For the attractive delta well at $x=0$ set
$\beta=m\alpha/(\hbar^2k)$.
The matching conditions give amplitudes
$r=\ii\beta/(1-\ii\beta)$ and
$t=1/(1-\ii\beta)$. Left--right symmetry gives
\begin{equation}
\boxed{S_{\delta}=
\frac1{1-\ii\beta}
\begin{pmatrix}\ii\beta&1\\1&\ii\beta\end{pmatrix}}.
\end{equation}
Indeed $|r|^2+|t|^2=1$ and
$r^*t+t^*r=0$, so $S^\dagger S=I$.
For a finite symmetric well of half-width $a$,
write $k=\sqrt{2mE}/\hbar$,
$l=\sqrt{2m(E+V_0)}/\hbar$,
and
$D=\cos(2la)-\ii(k^2+l^2)\sin(2la)/(2kl)$.
Solving the four continuity equations yields
\begin{equation}
\boxed{t=\frac{\ee^{-2\ii ka}}{D},\qquad
r=\frac{\ii(l^2-k^2)\sin(2la)\ee^{-2\ii ka}}{2klD}},
\qquad
\boxed{S_{\rm well}=\begin{pmatrix}r&t\\t&r\end{pmatrix}}.
\end{equation}
The probabilities $|t|^2$ and $|r|^2$
satisfy $T+R=1$; parity symmetry makes
the left and right reflection/transmission
probabilities equal.

\problem{2.54}{Transfer matrices, composition, and double deltas}
A transfer matrix relates right-hand amplitudes
to left-hand amplitudes:
\begin{equation}
\begin{pmatrix}F\\G\end{pmatrix}
=M\begin{pmatrix}A\\B\end{pmatrix}.
\end{equation}
Starting with
$B=S_{11}A+S_{12}G$ and
$F=S_{21}A+S_{22}G$, solve for $G$ and $F$:
\begin{equation}
\boxed{M=\frac1{S_{12}}
\begin{pmatrix}-\det S&S_{22}\\-S_{11}&1\end{pmatrix}}.
\end{equation}
Conversely, solving $G=M_{21}A+M_{22}B$
for $B$, then substituting into $F$ gives
\begin{equation}
\boxed{S=\frac1{M_{22}}
\begin{pmatrix}-M_{21}&1\\\det M&M_{12}\end{pmatrix}}.
\end{equation}
Hence
\begin{equation}
R_l=\left|\frac{M_{21}}{M_{22}}\right|^2,\quad
T_l=\left|\frac{\det M}{M_{22}}\right|^2,\quad
R_r=\left|\frac{M_{12}}{M_{22}}\right|^2,\quad
T_r=\frac1{|M_{22}|^2}.
\end{equation}
For a real localized potential $\det M=1$.
Two isolated potentials traversed from left to
right act consecutively, so
$\boxed{M_{\rm total}=M_2M_1}$.
For $V=-\alpha\delta(x-x_0)$, the
jump condition gives
\begin{equation}
\boxed{M(x_0)=
\begin{pmatrix}
1+\ii\beta&\ii\beta\ee^{-2\ii kx_0}\\
-\ii\beta\ee^{2\ii kx_0}&1-\ii\beta
\end{pmatrix}},
\qquad\beta=\frac{m\alpha}{\hbar^2k}.
\end{equation}
The two-delta transfer matrix is
$M(a)M(-a)$. Its lower-right element is
$M_{22}=(1-\ii\beta)^2+\beta^2\ee^{4\ii ka}$,
so
\begin{equation}
\boxed{T_{\delta\delta}=
\frac1{\left|(1-\ii\beta)^2+\beta^2\ee^{4\ii ka}\right|^2}},
\end{equation}
agreeing with Problem 2.28.
